  \documentclass{article}
  \usepackage[whole]{bxcjkjatype}

  \usepackage[numbers]{natbib}

  \usepackage{amssymb}
  % Language setting
  \usepackage[english]{babel}

  % Set page size and margins
  \usepackage[letterpaper,top=2cm,bottom=2cm,left=3cm,right=3cm,marginparwidth=1.75cm]{geometry}

  % Useful packages
  \usepackage{amsmath}
  \usepackage{graphicx}
  \usepackage[colorlinks=true, allcolors=blue]{hyperref}

  \usepackage[letterpaper,top=2cm,bottom=2cm,left=3cm,right=3cm]{geometry}

  % 段落の字下げと行間の設定
  \setlength{\parindent}{0em} % 全角幅の字下げ
  \renewcommand{\baselinestretch}{1.5} % 行間の調整

  \makeatother


  \begin{document}


% Title page
\begin{titlepage}
    \centering
    \vspace*{5cm} % タイトル位置の調整
    {\LARGE 科学技術と日本の将来\\[1.5cm]}
    {\Large \textbf{AIの安全な利活用システムMAGIの提案と用例}\\[4cm]}
    
    \vfill % 空白を埋める
    \raggedright
    \large 学校名: 都城工業高等専門学校\\
    学部・学科: 物質工学科\\
    学年: 専攻科１年生\\
    氏名: 福島武\\
    郵便番号：885-0011\\
    住所：宮崎県都城市下川東二丁目13号14番地-3\\
    電話番号:070-4104-1240\\
    E-mail:takeru.f.2004@gmail.com\\
    \vspace*{2cm}
\end{titlepage}


% Main content
\section{はじめに}

　「本文は誓って、手ずから人の手によって書かれている」

　この文言を論文の査読者が見た時、その査読者はその真偽を判断しなければならない。そのプロセスを如何にして行うのか、その審議がいま、様々な場所で行われている\cite{mbsj2023ai_science_forum}\cite{editage2024icmje_ai_guidelines}。そして、生成系AIの利用に関する利点とそこから生じる問題に関しても様々な議論が行われている\cite{cstp_ai_strategy}\cite{mext2024ai_innovation}。本文は、その議論に対する一つの解を提供したい。その方法は極シンプルで、あくまでも提出者の責任のもと生成系AIの利用を全く制限しないというものである。そのような強気なルールでも本当に価値の高い成果物を出そうと思えば自然と人間はどこか必要な場所で自分の手を介入させなければならないというのがこの主張の前提になっている。しかし、生成系AIは様々な雑務において革命的に有用であり、人間が行うべき本当に価値のある仕事に集中するための最高の道具なのである。

　本文は以下の主張によって構成される。
 
 　第一に、生成系AIによって生成されたものと人間が手ずから作ったものは区別がつかない。
  
  　第二に、生成系AIも人間も原理・理論から考えてほとんど同様に間違いを犯す。しかし、人間はその責任を負うことができる。
   
　第三に、2つ目の主張を受けて、著者が定義した新しい安全なAI利用システムである「MAGI」についてその概要と利用法を紹介する。
    
　そして、最後に以下の主張を述べ、本文を締めくくる。
    
　「本文は誓って、手ずから人の手によって書かれている。しかし、その真偽の判定は理論上、不可能である。いずれにしても本文についての責任の一切は著者自らが負うこととする」

　もちろん、AIの利用によって生じるリスクは画一的なものではなく、予測できていないものも多くある\cite{cstp_ai_strategy}。人間が責任を負うからと言って解決しない問題も多く顕在することも予測される。しかし、AIを正しく利用しようとしている多くの人々が不当にAIの利用を制限される未来は工学研究に携わっている著者の立場からもあまり良い未来とは思えない。ここで主張されるシステムはあくまでも工学研究者としてAIを便利に活用していきたい人間が考えたシステムであることは今一度強調しておきたい。

　本文を持って、日本におけるAIの活用が正しく、また安全な方向に向かうことを願う。


\section{生成系AIと人間の区別の不可能性}

　まず、直感的に考えてみたい。生成系AIにこうお願いしてみよう。 「ありがとうと書いて」こうして出てきたテキスト、「ありがとう」がAIが書いたものなのか、人間が書いたものなのか、その真偽を判定することは不可能である。この「ありがとう」に筆跡はない。今度はセンテンスを考えてみる。「本文は誓って、手ずから人の手によって書かれている」このセンテンスにおける「手ずから」にはある癖のようなものが感じられる。生成系AIもその性質上文章が長くなればなるほどある種の癖のようなものが露見してくる。実際に使ってみればそれを直感的に感じてもらえると思う。すなわち、そのようなAIに関する癖を判定することができれば、人間と生成系AIの区別が行えることになる。

　しかし、これは技術的に大きな矛盾を抱えている。そもそも、2024年12月19日現在、AIは一つではない。事前学習済みmodelという大きな枠組みでくくってもOpenAIのChatGPT\cite{openai_chatgpt}, Anthropicのclaude\cite{claude_ai} , metaのllama\cite{llama} ,xAIのGrok\cite{x_ai}など、様々なmodelがあり、さらにそのmodelを基に開発者たちはいかようにも”癖”を変容できる。「査読者に悟られないように著者の書き方を模したmodelを生成する」といった開発など、かなり容易にできる。それら全てのmodelについての癖を列挙するのはおそらく不可能である。

　さらに、実際に使用が判定できたとして、今度は人間が手ずから書いたものが「AIが書いたものだ」と誤認してしまう可能性が生じる。これは非常に皮肉な話だが、確実に起こり得る。なぜならば、先程「ありがとう」で示したように、「本文は誓って、手ずから人の手によって書かれている」と生成させることも可能だからである。

　以上、二点の理由により直感的にはAIと人間の区別を行う取り組みは途方もない堂々巡りによるリソースの奪い合いが生まれることが推測される。著者個人としてはそのような無益な運動に貴重な国力が割かれるべきではないと考えている。

　問題に対する世界の最先端の取り組みはどうなのか、技術的かつ社会的に俯瞰して眺めてみたい。

　現在、AIと人間の生成物を区別する「watermarking」と呼ばれる技術がある\cite{arxiv_241118479}。これは、AIの進展によって汎用的なAImodelに関する人間とAIの区別がつかなくなってきた問題への解を提供するために注目されている技術であり\cite{arxiv_241118479}、この手法はあらかじめテキスト生成の段階で人間とAIの生成の結果が識別可能なKeyを埋め込むことで事後的にAIの結果と人間のそれを区別させる技術である。2024年12月20日現在ではGoogle deepmindの提供する「SynthID」\cite{deepmind_synthid}という製品も提供されている。しかし、あらかじめkeyの埋め込みを行うことを前提にしたAI-modelでなければこの技術は活用できず、それを各国全てのAI-modelに適用するのは現実的に不可能である。

　すなわち、直感的にも技術的にも、AIと人間の生成の結果を区別することは現状、不可能であると言える。

\section{AIと人間の論理的な類似点と相違点について}

　論理的にAIと人間の類似点と相違点を考えたい。以降は様々な学問領域の知識を用いて、専門的な表現で論述を行うため、あらかじめその主張の要旨を以下に示しておく。

　まず第一に、AIと人間に論理的な差異は無い。

　すなわち、ある社会システムにおけるAIと人間の「構成要素としての役割」は全く等価である。

　しかし、人間は精度と速さで劣り、AIは人権と身体を持たず、また社会的責任を負うことができない。

　したがって、人間が社会的責任を負うことを前提にAIに「仕事」をさせるシステムが最適だと著者は主張する。そのようなシステムのことをMAGIと名付ける。

　以降はこの主張の各部における論拠を示す。

\subsection{AIと人間に論理的な差異は無い}

　ここでもまず、直感的な話をしたい。AIとは一言で言えばなんだろうか。AIは人間の知能の仕組みを構成論的に解き明かそうとする学問分野である\cite{matsuo2021}。人工知能（AI）はその時々で副産物として最先端の計算技術を生み出してきた\cite{matsuo2021}。すなわち、いま、ここで議論している所謂AIは「AIという学問分野の中で副産物として生まれた計算技術のうちの一つ」であり、他のその類の計算技術と明確に呼び名を分ける必要性があるといえる。以降は議論の中核となっている所謂AIのことを生成系AI（generativeAI,GenAI）と呼ぶこととする。「1.はじめに」では誤解を生まないよう、あらかじめ生成系AIと書いている。

　では、生成系AIとはなんなのか。生成系AIはニューラルネットワークであり、ニューラルネットワークは脳の構造を模して設計された計算アルゴリズムである\cite{matsuo2021}。ここで強調したいのは、AIとは人間の脳を模して作られた計算アルゴリズムであるということである。そして、「生成系AIを使うことの安全性」などという議論が挙がる程、「人間の脳を模した計算アルゴリズム」は「実際の人間の脳」と区別がつかなくなってきている。

　現在、ほとんどすべての社会システムを構成している人間という一つのニューラルネットワークも、言わずもがな”間違い”を犯す。その脳を模したニューラルネットワークが同様に間違いを犯すのは”直感的には”実に自然なことのように感じられる。では、AIは人間と同様に扱っても問題ないのだろうか。直感的には、なんだか「怖い感じ」がする。ここで、AIの利用とその結果生まれる恐怖についてのコンフリクトを正しく認知し、その恐怖を乗り越えるべく、以下の論述でこの議論の解を提供したい。

　論理的に反証の余地を可能な限り残さないように説明する。

　ラッセルに代表されるような数理哲学、グレアム・ハーマンに代表されるようなオブジェクト指向オントロジー\cite{harman2018}の概念を数学の集合の概念とその記述法を用いて定義し、「認識しうる現実を最も善く表現している論理空間」を数学的に定義した上で、その論理空間において情報を伝達する構成要素を主体Sと名づけ、定義する。結果的に、この論理空間において人間もAIも主体Sの集合に包含可能であり、その観点でAIと人間に差異は無いと言うことを説明する。

　まず、世界に存在する全ての「論理」すなわち「論理命題」を定義する。


\(\quad P \) を世界に存在する認識可能な全ての論理命題からなる有限個の要素を持つ集合として以下のように定義する。

\[\begin{aligned}P &= \{ p_1, p_2, \dots, p_n \mid p_i \text{ は認識可能な論理命題}, \; N,n \in \mathbb{N}, \; n \leq N ,i=1,2, \dots,n\}  &\quad \\&\begin{aligned}1. \; &\forall p_i \in P  \text{ は真または偽の値を持つ } (p_i \in \{\text{真}, \text{偽}\}).  
\\2. \; &\forall p_i \in P\text{に対する}
\; \Bar{p_i}  \text{（すなわち } p_i \text{ の否定命題）もまた } P \text{ に属し、真または偽の値を持つ}. 
\\3. \; &p_i \text{ が自己言及的命題であり、その真偽が定義不可能な場合、} p_i \text{ は } P \text{ に含まれない}. 
\\4. \; &\text{自己言及パラドクス（例: 「この命題は偽である」）により、真偽が判定不能な命題は、}\\   &\text{定義上、論理的に認識可能な命題の集合 } P \text{ の外部にあるものとする}.\end{aligned}\end{aligned}\]


　難しいことを書いているが、要は「全ての論理」を定義したと思ってもらえればいい。例えば、「りんごは赤い」などもその一つである。これは概ねの場合「真」という要素を持つ命題になる。例えば「りんごは青い」も同様である。そのような命題が世界には「有限個」存在していると定義できる。無限にあるようにも感じるかもしれないが、主体（生物や人間）の数は有限である場合、それらが認識しうる論理命題もやはり有限個になるからである。

　次に、例えば「りんご」という単語、対象そのものについて考えたい。これは論理命題ではないが、我々は確かにそれを「認識」している。眼の前に赤いヘタのついた果物を眼の前にしたとき、我々は「りんごだ」と思う。しかし、そもそもそれは「りんご」という名を冠する以前から存在しているもの（対象）そのものなのである。このような「対象そのもの」のことをオブジェクトと呼ぶことにする。

\[\text{主体の認識にかかわらず世界に存在する全ての対象をオブジェクトと呼ぶ。}\]
\[\text{それらを集めた集合を } O \text{ と定義する。}\]
\[O = \{ o_1, o_2, \dots, o_n \mid o_i \text{は認識可能なオブジェクトそのもの},N,n\in \mathbb{N},n \leqq N, j=1,2,\dots,n \}

　ここで各オブジェクト\(o_j \in O\)は、それを説明する論理命題の部分集合によって定義できる。つまり、任意の\(o_j\)に対して、

\[o_j \leftrightarrow P_{j} \subseteq P, \quad |P_{j}| \leq N.\]

が成り立つ。さらに、各\(p_i\in P_i\) もそれを表現する時、何かしらのオブジェクトの部分集合になっている。


\[p_i \leftrightarrow O_{p_i} \subseteq O, \quad |O_{p_i}| \leq N\]

　例えば、「りんごは赤い」という命題は「りんご」というオブジェクトと「赤い」というオブジェクトに分割可能である。もちろん、「り」「ん」「ご」「は」....などにも分割可能である。

　それらのオブジェクトはさらに別の論理命題の部分集合によって定義される。このあるオブジェクトを定義するための論理命題を定義するためのオブジェクトを定義するための…という連鎖は無限に続く。
\[o \to P_{o} \to O_{p} \to P_{o'} \to O_{p'} \to \cdots\]
　ここで、この連鎖の中における論理命題の部分集合の不一致性について着目したい。「りんご」というオブジェクトを定義する時、「りんごは（一般的に）赤い」という命題が選択されたとする。次に、「赤い」するため「血は赤い」という命題が定義として選択されたとする。今度は「血」というオブジェクトについて「血が出ると絆創膏によって傷を塞ぐのが一般的だ」と言う命題が選択されたとする。ここで、一度俯瞰してみると、りんごというオブジェクトの定義の中には「絆創膏」というオブジェクトが含まれてしまったことになる。このような意味の揺れ動き、初め定義したはずの論理命題とオブジェクトの要素の増加は止まることはない。集合\(O\)も\(P\)も有限個の要素を持つことから、いずれ、一つのオブジェクト、ここではりんごを説明するための定義の連鎖は全ての論理命題とオブジェクトの集合を含んでしまうことになる。これは、「どのオブジェクトもいかようにも定義可能であり、どの論理命題も如何様にも表現可能である」ということを示している。その部分集合も有限個の要素を持つ。このような意味とその表現に関する連鎖はフランス現代思想におけるラカンの言葉を借りるとシニフィアン（の）連鎖\cite{kawasaki2007}であると言える。精神分析と強く結びつくフランス現代思想のシニフィアン連鎖と論理学から出発した数学的な「認識論」が結びついたのは、「真なる思考の総体が世界の像であ」（三・〇一）り、「思考は命題において知覚可能な形で表され」（三・一）、「我々は可能な状況を射影するものとして命題という知覚可能な記号（音声記号、文字記号、等々）を用いる。命題に対してその意味を考えること、それはすなわち射影方法を考えることにほかならない」（三・一一）こと\cite{wittgenstein1922}に関係していると考えられる。すなわち、考えることも伝達も「主に」言葉によって行われることに関連しているからだと言える。

　ここで一つ、目的を見失わないためにこの論述で主張したい内容を今一度確認する。ここで主張したいのは、この論理空間において「言葉を話す人間とAIはシステムの構成要素として論理上、同じ集合に含まれ、論理空間上はAIと人間に差異は無い」ということである。したがって、以降はAIと人間をある一つの集合\(S\)と名付け、AIと人間を同一の主体集合\(S\)でくくっても矛盾性をはらまないことを確かめる。

\[ \text{主体}  S \text{は、「認識を行うもの一般」の全集合のことを指す。}\]
　各主体が認知可能な論理命題の全集合\(P\)  と認知可能なオブジェクトの全集合 \(O\)  の部分集合に関連付けられる集合として定義する。
\[S = \{ s_{ij} \mid s_{ij} = ( P_{i}, O_{j}) \}, \quad i \in \{1, 2, \dots, n\},\quad j \in \{1, 2, \dots, m\} \quad n,m \in \mathbb{N}\]
ここで、
\[P_{i} \subseteq P \quad (\text{主体 } s_i \text{ に関連する論理命題の部分集合})\]

\[O_{j} \subseteq O \quad (\text{主体 } s_j \text{ に関連するオブジェクトの部分集合})\]

　すなわち、ある主体はすでに認識したことのある有限個の要素を持つオブジェクト集合とそれを定義する有限個の要素を持つ論理命題の集合とで「論理空間上は」定義できる。そして、2つ以上の主体間では情報伝達が行われ、以下のように定義する。

 
\[T(s_{ij}, s_{lm}) = \begin{cases} (p_s, o_s) &  o _s\notin O_{s_{lm}}\text{かつ} p_s \notin P_{s_{lm}} \text{ ならば、}s_{lm}\text{内で一定の確率 } q \text{ で}\text{生成} \\(\bar p_s, o_s) &  o _s\in O_{s_lm} \text{ かつ } p_s \in P_{s_{lm}} \text{ ならば、}s_j\text{内で一定の確率 } q \text{ で真偽反転}\end{cases}\]

　ここで右辺の\((p,o)\)は\(s_{lm}\)に生成あるいは改変された論理命題とオブジェクトの部分集合の事を指す。

　受け手主体が提供されたオブジェクトあるいは論理命題を知らなかった場合、一定の確率 \(p\) で伝達されたオブジェクトとその論理命題が受け手主体の論理空間に生成される。

　伝達された内容を知っているオブジェクトと論理命題だった場合、一定の確率 \(p\) でその論理命題の真偽が反転する。確率\(p\)は固体固有のいくつかの関数の集合を指す。これは生活の中、または相手との関係性の変化で変わるものとするが、本文では厳密な定義を避ける。


　情報伝達は論理空間上で、\(Q \neq 0\)の確率\(Q\)で必ず起こるものとする。これは、我々が生活を送る上で情報伝達が必ず起こることからそう定義する。

　そして、全ての主体の持つオブジェクトの集合とその論理命題の集合はシニフィアン連鎖の理論に照らすと、オブジェクトの全集合\(O\)と論理命題の全集合\(P\)を用いて「定義可能」であり、その全集合には初め定義したように偽となる論理命題も含まれる。

　すなわち、\(Q \neq 0,q\neq 0\)である限り、主体には嘘が伝達される可能性は0ではない。

　これは、主体の確率\(q\)の問題ではなく、そもそも、シニフィアン連鎖のただ中にいる論理空間の問題であり、主体は伝達を迫られる（\(Q \neq 0\)）であるからだと言える。

　したがって、論理空間上で人間とAIを主体Sで定義可能であるならば、人間であろうがAIであろうが、間違い（ハルシネーション）が起こる可能性は0ではない。このことから、あくまで論理空間上は人間とAIに機能上の差異は無く、その情報の真偽を行う確率\(q\)の関数を厳正に審査してシステムでの利用を判断するべきである。

　しかし、ここで社会的な問題を考えたい。AIには、まだ、というべきか、人権が無い。そのハルシネーションで人権を持つ主体に損害を与えた時、その損害を誰が保障するのだろうか。この矛盾を解決するために、あくまでも人間がその損害の保障に関する責任の一切を負うべきであると主張する。

　以上の論述を以下にまとめる。

　AIも人間も同様に間違いを犯す。

　しかし、AIは責任を負うことができない。

　したがって、人間が責任を追う必要がある。

　逆を言えば、人間が責任を負う前提であればAIを利用しても問題がない。

　ここまでの主張を受けて、このような人間が責任の一切を追うことを前提に生成系AIやその他のソフトウェアアプリケーションを使うシステムのことを著者はMAGIと呼称し、以下にその定義を示す。

\section{MAGIについて}
　本章は著者がBlogに書いている内容を改変したものであるためネット検索にヒットする可能性があるが断じて著者自らの主張である。

　AGIという言葉が巷では繁茂している。変な表現だが、定義があやふやなまま、言葉だけが滞留して恐怖と不安が煽られるツイートなどは決して珍しくない。 AGIは「Artifitial general interigence」の略で、実に様々な定義が専門感の中でも滞留しているのが垣間見える。\cite{morris2024}

　ここで僭越ながら著者も繁茂しているAGIの定義の一つを提案したいと考えている。 提案と言っても単に、今後論述を行っていくうえで生活でAGIという単語を使いやすくしたいから立場をはっきりさせたいだけなので、これを用いて揚げ足を取ろうなどとは考えていないしこれを読んでいる人もそう思ってほしくない。

　さて、既存のLLMのようなブラックボックスなプログラム群のことをここで「Black code」と称する。抽象的に「原理上、ブラックボックスなもの」全てをBlack codeと定義する。同様に、既存の自然科学の知見、プログラム言語、自然言語など「原理が記述可能なもの」をWhite codeと呼びたい。 そして、AGIは、AIがこのBlack codeに指示を出し、White codeを生成し続け、加速膨張し続けるシステムのことを指したいと考えている。Generalと入っているので大きくて統括的なシステムが想定されていることが多いように見受けるが、あえて、「自分でCLIを用いてフォルダやファイル中のコードを生成し続け、目的に応えていくシステム」のことを一般性を持ったGeneralなinterigenceだと解釈する。 ちまたには「Agent」と呼ばれる製品群\cite{jjsai2011}もあるが、AGIの「CLIを用いてフォルダやファイルを自立生成できる」（アプリケーションを自ら拡張できる）という点が大きな違いになる。つまり、Agentの行き着く先はAGIだと著者は考える。

　ここでよく考えたいのは、このAGIの倫理性である。倫理というのは何もスピリチュアルに言っているわけではなく、学会等で硬く守られているあくまでも”科学的な”文脈での倫理性を一度考えねばならないと思う。何故か。「倫理は人間の活動（ここでは行為全般のことを指すことにします）の規範を適切に定めるものであり、AGIは人間と同様に”活動”が行えるようになるから」である。 「WebのAGIを作りました。自律的に勝手に膨張して、著作権侵害、風紀違反しまくり」じゃ困る。そこには何か倫理を規定しなければならず、それを何百年、何千年も考えてきた科学がこの世界にはどうやらある。著者は今丁度倫理を習っていて、大きく功利主義、義務論、徳倫理学がある\cite{muramatsu2017}という認識を持っているが、AGIにも徳、つまり正解となる人物像、AI像を定めたうえで自立膨張、変化をしてもらわないといけなくなると著者は考える。しかし、ここで大きな矛盾が生じる。ヒトラーの倫理性を徳としてAGIが自立駆動したらどうなるだろうか。数理とコンピュータはすでに最低限の性能のものは誰でも使えるから規制も難しそうだ。 つまり、AGIはおそらく限りなく果てしのない善と悪のいたちごっこを生んでしまうというのが著者の考えである。

　しかし、往々にしてAIというのは便利なものだ。AGIのセキュリティや国家の対立等は工学者の著者は一度無視させてほしい。工学的に利用できるような正しいAIの使い方があるはずだ。そこで、MAGIという新しいシステムを提案したい。 正直なことを言うと、卒業研究の中で自分の提案するシステムのことを命名するときにパッとエヴァンゲリオンのコンピュータシステムであるMAGIシステムを思い出して名付けたという経緯がある。ほら、あの科学者の女の人のお母さんのやつ。「貴方は最後の最後まで母ではなく、女だったのよ」みたいな決め台詞がある。 失敬、MAGIはもともとキリスト教の聖書の中に出てくる三賢者のことを指す。\cite{takashina2011}AGIにおけるAI、Black code、White codeの三すくみをそれぞれ賢者と抽象化したうえでそのAIの部分を人間に置き換えてAGIを作ろうというのが著者の提案である。 むしろこの概念を説明するために先ほどAGIを定義したというのが本音である。

　MAGIにはいろんなMAGIがある。WebMAGI(Webを勝手に巡回して調べてくれるMAGI、Web Agentと呼ばれる仕組みを拡張)やRAGMAGI（与えたデータから抽出して回答してくれるLLMであるRAGの拡張）などなど、既存の全てのシステムはMAGIによって置き換えられるというのが僕の考えである。 人間が必要に応じてBlack codeに指示を出してWhite codeを生成し、そのWhite codeを人間が監査し、また必要があればBlack codeに指示を投げる、ただそれだけのシステムである。この人間のシステムに対する認知と関与がとても重要であると考える。その論拠は「2.AIと人間の論理的な類似点と相違点について」で説明したとおりである。このような、人の手によって監査・改善されるシステムであるMAGIはManus Artifitial general interigence.と呼ぶことにする。Manusはラテン語で「手」と言う意味があるそうである\cite{cambridge_manus}。そして、Manusはキリスト教の文脈で「Hand of god」、すなわち神の手と訳され神が世の中の出来事に介入することを示すために用いられる。しかし、ここで一つ「神」という言葉のシニフィアンを連鎖させたい。ここで述べられる神はキリスト教的な、唯物神的な、西洋中心主義的な所謂「神」ではない。もっと汎神論的、多神論的、アニミズム的な\cite{ishizuka2020}「神」の手（Manus）として定義したい。僕らの世界に神（唯物神）などいない。いや、いてはいけない、と言ったほうがいい。いたとしたらそっぽを向いて他の宇宙でも作っていてほしいものである。この宇宙に居る主体は文部科学大臣だろうが、20歳の田舎のただの高専生だろうが、提出前にこの論文を査読してくれている多くの都城高専の先生方であろうが、あるいはただの石であろうが宇宙の中では等価である。それは、もちろん、AIであってもそうで、ここで指す「神」は「＝人間」と訳さないで欲しい。あくまでも、考えうる全ての考慮のうち、この役割を担うのが今は偶然「人間であった」というだけである。いずれ、AIについての理解が深まっていけばポツポツとその役割の一部を機械に代替してもらうときもやってくると著者は考える。まずは、AIがアニミズム的な意味で人類と同類なのだ、ということを多くの人間が理解することが重要なのではないかと著者は考える。そのためにもまずは人間を監査に据えたMAGIシステムはかなり有意義な意味を与えるのではないだろうか。

　MAGIの定義について、最後になるが、このMAGIシステムはテクノロジーと人間との調和をとるためのかなり善いシステムだと直観している。著者の宇宙で一番尊敬している落合陽一先生\cite{scholar_citations}はかなり昔から「デジタルネイチャー」\cite{ochiai2018}という概念を提唱しており、デジタルネイチャーは話せば長くなるが「コンピュータと既存の自然により形成された新しい自然」というのが簡単な定義だといえる\cite{ochiai2018}。実用的な視点を言うと「機械が”偶然性”を獲得した状況」のことを指していると著者は認識している。偶然性というのは哲学者である千葉雅也氏から引用している\cite{utcp_tokyo}。これまでの機械はテキトーやほんとうの意味でのランダムは創れなかったのに対し、Black codeのLLMは簡単にテキトーな対象を生成してくれる。これはある種、既存の自然に照らすとおかしなことなのである。 「自然は分からない」だからわかろうとする、それが科学の営みである。自然とはBlack codeらしい。人間の作ったものは、理解ったものから作っているはずだから、Black codeになりえない。それがこれまでの一般認識だった。しかし、いまの機械はどうだろうか。原理的に理解できないものが出てきてしまった。「人間の作ったものの一部」もわからない、でもわかろうとしないといけない。わからないまま使うなんてできない。この自然と人間のやりとりの一般性に照らして生じる矛盾を解消するために、「計算機も含めて新しい自然だと認識してもいいじゃないか」と言っているものと著者は認識している。つまり、MAGIシステムにおけるBlack codeはデジタルネイチャーであり、White codeはその自然を解明して人間が認知を行い理解ったもの、そして最後に人間そのもの、と新しい宇宙の構造を定義し直せる。 「デジタルネイチャー」はあくまでも事実を言葉として定義しただけに過ぎない。著者は、MAGIシステムをもってして初めてデジタルネイチャーと人間との関係性とその在り方を定義したい。そして、それをポストデジタルネイチャーの基盤的な考えにしたいと勝手ながら考えている。その概念に照らすとMAGIシステムはデジタルネイチャーでしかない。人間も機械も、宇宙の構成要素として等価である。しかし、やはり僕らと機械は環世界\cite{hatagami2017}が違う。同じ宇宙の構成要素である機械と人間を環世界が違うという理由で区別している。その古典的な自然感を僕はあえて、ポストデジタルネイチャーと名付けたい。ポストであるにもかかわらず、その概念はデジタルネイチャーの内側に包含されている。つまり、ポストデジタルネイチャーの概念はデジタルネイチャーの概念に3つの塗り絵を行ったと考えるといいかもしれない。あくまで、人類が自然を解明しようとする時、自然と人類は対比されてきた\cite{seki1997rethinking}。もちろん、厳密には僕らも自然の一部なのだけど。 著者の中でポストデジタルネイチャーはそういう二重性を持った概念として直観されている。しかし、何度も言うように機械も人間もその存在論的な価値は等価である。ポストデジタルネイチャーは、デジタルネイチャーの次世代的な考え方でありながら、その意味はデジタルネイチャーに完全に内包されている。

\section{MAGIの利用}

　このような考えのもとシステムを運用すれば、これまでと同じ査読のシステムにあくまでも提出者の責任のもと、提出者は生成系AIの使用を行うことが出来る。この通常の、「生成系AIが出る前まで行われていた規則」を用いることが非常に重要だと考える。この点からMAGIシステムを眺めると、既存の「（已然、間違いを犯す可能性のある）人間が責任を担保することで」運営されている全てのシステムに、MAGIは置換可能であると推測される。
　しかし、本文の主張は、「わざわざ生成系AIの使用を制限する必要は実のところあまりない」とも言い換えられる。例えば、あるタスクにおける一定以上のスキルを持っている人間は、わざわざ生成系AIを使わずとも自分でそのタスクを行ったほうが、監査の時間も考慮すると結果として早く済むことも少なくない。例えば、この論文を査読してくれた都城高専のある国語の教員の方は、生成系AIを用いた簡単な実験を授業を通して行っている。内容は簡単で、エントリーシートを模した文章をいくつか生成させ、その中に教員の書いたものを織り交ぜ、そのいくつかの文章のうちから「良い」と思われるものを選ばせるというもので、先生によればやはり人間が書いたものが一番選ばれる割合が高いらしい。もちろん、生成系AIに指示を出せばその精度は次第に上がるのだろうが、いまのところ、人間が成果物からその人間そのものを評価するとき、その評価はまだ人間が生成物を作ったときのほうが高いようである。すなわち、わざわざ生成系AIの利用を制限しなくても、本当に人間として味があってい良い成果物を出そうと思えば自然人間は自分の手を介入させるのではないかというのも著者の考えであり、それに必要な補助的な雑務、調査やちょっとした言い換え、反論の探索などに生成系AIは有用だと著者は考える。

　いずれにしても生成系AIは使用を制限されるべきものではないと考える。

\section{おわりに}

　以上の論述を再度以下にまとめる。

　第一に、生成系AIによって生成されたものと人間が手ずから作ったものは区別がつかない。

　第二に、生成系AIも人間も原理・理論から考えてほとんど同様に間違いを犯す。しかし、人間はその責任を負うことができる。

　第三に、2つ目の主張を受けて、AIの生成結果を人間の責任において提供するシステムである「MAGI」についてその概要を紹介した。

　最後に、「本文は誓って、手ずから人の手によって書かれている。しかし、その審議の判定は理論上、不可能である。いずれにしても本文についての責任の一切は著者自らが負うこととする」


\bibliographystyle{unsrt}
\bibliography{sample}

\end{document}
